\chapter{CONCLUSÃO}

A comparabilidade entre métodos automáticos convencionais e gravimétricos para \MPdez{} e \MPdoiscinco{} varia com o sistema completo de medição e suas condições de operação. Aquecimento, semivoláteis, umidade, propriedades do aerossol, concentração e tratamento dos dados explicam parte das diferenças. O sinal do desvio não é fixo por tecnologia, e os artefatos do comparador gravimétrico também precisam ser considerados.

Correções locais podem melhorar o desempenho quando calibração e validação são distinguidas e a aplicação permanece no domínio avaliado. O conjunto não sustenta um fator universal nem uma média global de desvio. Para monitoramento e análise temporal, é necessário documentar fração, configuração, referência, integração, controle de qualidade e mudanças de método. As conclusões permanecem condicionadas à disponibilidade documental, à heterogeneidade e às limitações da revisão por pesquisador único.
